\section{Finite-moment models for the lamp source}\label{sec:power-moments}
We use the six source unitaries $p,q,t,d,x,z$ and the $27$ words of~\cite[Definition~7.1]{DM-game}. Set $s=[p,q]$, $B=psp^{-1}$, $t_i=p^{i-1}tp^{1-i}$, $J=(xz)^2$, and normalize the residuals by $R_w=(w-I)/2$, together with $R_J=(J+I)/2$. Every residual has norm at most one on unitary assignments.

\begin{theorem}[All fixed finite moments]\label{thm:power-moments}
For every fixed real $\nu>0$ and every sufficiently large integer $N$ there are flat unitary assignments of this same source with $x^2=z^2=I$, $J=-I$,
\[
 \max_w\tau_{D_N}(|R_w|^{2\nu})\le C_\nu N^{-\nu},\qquad
 \log D_N=3N+O_\nu(\sqrt N),\qquad R_J=0.
\]
The presentation does not depend on $\nu$. No uniform estimate as $\nu\to\infty$ is claimed.
\end{theorem}

\subsection{Character moments and Clifford lamps}
Fix $\nu\ge1$, $k=\lceil\nu\rceil$, $m=\lceil\log N\rceil$ and $F=\mathbb F_{2^m}$. Choose distinct $\alpha_1,\ldots,\alpha_N\in F$ and a binary basis, and put
\[
 \zeta_j=(1,\alpha_j,\ldots,\alpha_j^{2k-1})\in\mathbb F_2\oplus F^{2k-1}\cong\mathbb F_2^{n_0},
 \quad n_0=1+(2k-1)m.
\]
Pad into $W=\mathbb F_2^n$ for any $n\ge n_0$. Define
\[
 v_u=N^{-1/2}\bigl((-1)^{\langle\zeta_j,u\rangle}\bigr)_{j=1}^N,
 \qquad c(u)=N^{-1}\sum_{j=1}^N(-1)^{\langle\zeta_j,u\rangle}.
\]
Then $\langle v_u,v_w\rangle=c(u+w)$ and, with $b_k=(2k-1)!!$,
\begin{equation}\label{eq:power-character}
 \frac1{2^n-1}\sum_{u\ne0}|c(u)|^{2\nu}\le b_k^{\nu/k}N^{-\nu}.
\end{equation}
To prove this, any binary dependence of at most $2k$ distinct columns has even length, by its first coordinate. For the dependent column set $S$, its field power sums of orders $0,\ldots,|S|-1$, with $|S|\le2k$, then vanish. The Vandermonde matrix on the distinct $\alpha_j$ would annihilate the all-ones vector, a contradiction. Thus the signs are $2k$-wise independent and unbiased for uniform $u$. In their $2k$-th moment only even multiplicities survive; pairing positions bounds their number by $b_kN^k$. Divide by $N^{2k}$. Removing $u=0$ decreases the average since $c(0)=1$ and $|c|\le1$, and concavity of the power $\nu/k$ proves \eqref{eq:power-character}.

Use full exterior Fock space $\mathcal H_0=\Lambda(\C^{3N})$, of dimension $2^{3N}$, with Majoranas satisfying $\{\gamma(f),\gamma(g)\}=2\langle f,g\rangle I$ on real vectors. In three orthogonal copies set
\[
 A_u=\gamma(v_u\otimes e_1),\quad B_u=\gamma(v_u\otimes e_2),\quad C_u=\gamma(v_u\otimes e_3),
 \qquad X_u=iA_uB_u,\quad Z_u=iA_uC_u.
\]
These are reflections, $X_uZ_u=B_uC_u=-Z_uX_u$ and $(X_uZ_u)^2=-I$.
The Clifford construction of overlapping qubits is due to Chao, Reichardt, Sutherland and Vidick~\cite[Theorem~3.1]{CRSV17}. We use its three-family Clifford construction; the fixed finite moments and simultaneous affine covariance needed here are established below. For $c=c(u+w)$, CAR expansion gives
\[
 X_uZ_w-Z_wX_u=2cB_uC_w,\qquad
 X_uX_w-X_wX_u=-2c(A_wA_u+B_wB_u-2cI).
\]
Since $B_uC_w$ is unitary, the absolute normalized mixed group residual is exactly $|c|I$. For the Majorana reflections $A_u,A_w$ (and likewise $B_u,B_w$), the identity $UV+VU=2cI$ gives $(UV-cI)^*(UV-cI)=(1-c^2)I$. Put $P=A_uA_w$ and $Q=B_uB_w$. They commute, $X_uX_w=PQ$, and the same-type group commutator is $(PQ)^2$. Since $P+P^*=Q+Q^*=2cI$, write $P=cI+i\sqrt{1-c^2}R$ and $Q=cI+i\sqrt{1-c^2}S$ for commuting reflections when $|c|<1$. The scalar joint-spectrum values give $\|(PQ)^2-I\|/2\le2|c|\sqrt{1-c^2}\le2|c|$. The endpoint cases follow directly, and the identical calculation applies to $Z$.

Let $D_u=\operatorname{diag}((-1)^{\langle\zeta_j,u\rangle})$ and $U_u=\Lambda(D_u\otimes I_3)$. Full second quantization is a genuine representation, so $U_uU_w=U_{u+w}$ and $U_uX_wU_u^*=X_{u+w}$, likewise for $Z$; there is no projective sign ambiguity.
On $\ell^2(G)\otimes\mathcal H_0$, $G=\operatorname{GL}(n,2)$, put
\[
 L_h(\delta_g\otimes\psi)=\delta_{hg}\otimes\psi,\qquad
 T_u(\delta_g\otimes\psi)=\delta_g\otimes U_{g^{-1}u}\psi,
 \qquad \widetilde X_w|_g=X_{g^{-1}w},\quad\widetilde Z_w|_g=Z_{g^{-1}w}.
\]
Direct substitution proves $L_hT_uL_h^*=T_{hu}$, $T_uT_w=T_{u+w}$ and
$L_h\widetilde X_wL_h^*=\widetilde X_{hw}$, $T_u\widetilde X_wT_u^*=\widetilde X_{u+w}$, with the same identities for $Z$.
Thus this is an exact affine representation. Every linear action commutes with the zero lamps. Its logarithmic dimension is at most $n^2+3N$.

Finite-field independence and small-bias spaces have established precedents in Naor--Naor~\cite{NN93} and Alon--Goldreich--H\aa stad--Peralta~\cite{AGHP92}. The Vandermonde argument and even-moment pairing are standard tools; the additional requirements here are simultaneous affine covariance, the literal source certificate and the collective clock. The finite exhaustive search in the small-bias side bound below is not claimed to be a new efficient derandomization.

\subsection{The collective clock in the affine representation}
Here we give the geometric lift explicitly. Its permutation regions and spreading method are those of~\cite[Propositions~6.1 and~6.5, Lemma~6.4]{DM-H3} and~\cite[Proposition~13.9]{DM-precision}.
Take $M\ge48$, $n=4M+4$, $\ell=2M+2$ and right-coordinate cycles $c_r=(1\ r\ r-1\ \cdots\ 2)$. At $0\le i,j<M$ use
\[
 A_{i,j}=c_{\ell-i-j},\quad B_{i,j}=c_{2\ell-j}\ (j<M-1),\quad B_{i,M-1}=D_i,
\]
\[
 D_0=c_{2\ell-M+1},\qquad D_{i+1}=c_{\ell-i-(M-1)}^{-1}(1\ 2)D_ic_{\ell-i}.
\]
Make the preliminary exact moves $(1\ M+4)$ at the $a$-link $(M-1,M-1)$ and $(5\ M+5)$ at the $b$-link $(0,M-2)$, multiplying on the right-action left. The only exceptional plaquettes have labels
\[
 P_B(q_B)=\prod_{r=5}^{M+3}(r\ r+M),\quad P_T(q_T)=(5\ M+5),\quad
 q_B=(M-1,M-1),\ q_T=(0,M-2).
\]
The bulk region contains $q_B$, the column $(M-1,j)$ for $1\le j\le M-2$, and $0\le i\le M-3$, $0\le j\le M-2$, $i+j\le M-3$.
The top region contains $q_T$ and $i,j\ge0$, $i+j\le M-3$. Their areas are
\[
 A_B=M(M-1)/2,\qquad A_T=1+(M-1)(M-2)/2.
\]
The $r$th transported bulk transposition at a vertex $(i,j)$ has endpoints
\[
 \{r+M-1-i-j,r+2M-1-j\}\ (i\le M-2),\qquad
 \{r+M-1-j,r+2M-1-j\}\ (i=M-1).
\]
Take their product for $P_B(i,j)$; the top label is $\{M+3-i-j,2M+3-j\}$. The cited covariance identities restrict to this bulk region: labels are disjoint, base labels avoid $1,2$, and initial labels of active interior edges avoid $2,3,4$. Only those edges carry potentials; the source is attached by its bottom edge. These are finite permutation identities, so they can be checked directly from the displayed cycles and labels.

Use the inverse coordinate map to pass from right coordinates to left linear actions, and represent through $L$ above, denoting the representation by $\Pi$. Set $Q_R(v)=(I-\Pi(P_R(v)))/2$. Then $Q_R$ is a projection, $e^{i\pi Q_R}=\Pi(P_R)$, covariance survives and commuting labels give commuting projections. A label fixing $2,3,4$ commutes with translations on those coordinates and the transvection $E_{3\leftarrow2}$; its represented projection and all its exponentials do too. Every linear action fixes configuration zero, so all these corrections also commute with both zero lamps. This verifies the protected-affine operations rather than assuming that an arbitrary representation acts on named qubit factors.

Solve $\operatorname{curl}\varphi_R=\pi/A_R-\pi1_{q_R}$ with zero boundary potentials. The required curls sum to zero; eliminating leaves of the connected dual spanning tree assigns each parent edge to meet its leaf curl, and leaves the root equation automatically. Use the corrected links
$F_a(v)=\prod_R e^{i\varphi_R(a\text{-edge})Q_R(v)}\Pi(A_v)$ and similarly $F_b$. Define
\[
 p|c,\psi\rangle=|c-e_1,F_a(c-e_1)\psi\rangle,\qquad
 q|c,\psi\rangle=|c-e_2,F_b(c-e_2)\psi\rangle.
\]
Covariance brings each correction in a plaquette to its base. Commuting projections add their angles; at a source $e^{i\pi Q_R}$ cancels the original defect. The remaining holonomy is
\[
 e^{i1_B\pi Q_B/A_B}e^{i1_T\pi Q_T/A_T},\qquad
 \|[p,q]-s_0\|\le h_M:=\pi/A_B+\pi/A_T\le8\pi/M^2,
\]
where $s_0=\Pi((1\ 2))$. Put $t=T_{e_1}$, $d=L_{E_{2\leftarrow1}}$, $x=\widetilde X_0$, $z=\widetilde Z_0$.
The protected identities give
\[
 ptp^*=qtq^*=T_{e_2},\quad pT_{e_2}p^*=T_{e_3},\quad
 pdp^*=qdq^*=L_{E_{3\leftarrow2}}.
\]
Also $ps_0p^*=qs_0q^*=B_0=\Pi((2\ 3))$ and $p^2s_0p^{-2}=\Pi((3\ 4))$, so the independent-$s_0$ Houghton words are exact. The point swaps give $B_0s_0(pdp^*)s_0B_0=d$, and similarly with $q$, while $s_0p,s_0q$ commute with $t$.
The genuine affine identities give $t^2=I$, commuting coordinate translations,
$dT_{e_1}d^{-1}=T_{e_1}T_{e_2}$ and $[d,T_{e_2}]=[d,T_{e_3}]=I$. Every base/correction and $d$ commute with $x,z$; their onsite relations and $J=-I$ are exact.

For completeness, the entire $27$-word certificate has the following normalized bounds, with $h=h_M$, $K=\max_{u\ne0}|c(u)|$. Rows containing $s$ use its actual value $[p,q]$.
\begin{center}\small\setlength{\tabcolsep}{3pt}
\begin{tabular}{rll@{\qquad}rll}
\toprule
row & word & bound & row & word & bound\\\midrule
1 &$s^2$ &$h$& 2 &$(spsp^{-1})^3$ &$3h$\\
3 &$[s,p^2sp^{-2}]$ &$2h$&4 &$psp^{-1}qs^{-1}q^{-1}$ &$h$\\
5 &$t^2$ &$0$&6 &$[t,sp]$ &$h$\\
7 &$[t,sq]$ &$h$&8 &$[d,Bsp]$ &$2h$\\
9 &$[d,Bsq]$ &$2h$&10 &$[t_1,t_2]$ &$0$\\
11 &$dt_1d^{-1}(t_1t_2)^{-1}$ &$0$&12 &$[d,t_2]$ &$0$\\
13 &$[d,t_3]$ &$0$&14--19 &$[p,x],[p,z],[q,x],[q,z],[d,x],[d,z]$ &$0$\\
20--22 &$x^2,z^2,(xz)^4$ &$0$&23 &$[x,txt^{-1}]$ &$\le\min(1,2K)$\\
24 &$[x,tzt^{-1}]$ &$K$&25 &$[z,txt^{-1}]$ &$K$\\
26 &$[z,tzt^{-1}]$ &$\le\min(1,2K)$&27 &$[J,t]$ &$0$\\\bottomrule
\end{tabular}
\end{center}
The first four words have $2,6,4,2$ occurrences of $s^{\pm1}$; rows $6,7,8,9$ have $2,2,4,4$, after expanding $B$. Replace each by $s_0^{\pm1}$ and telescope, then divide by two. This proves every host bound. The zero rows follow from the displayed exact affine and onsite identities. For the four lamp rows, $t$ sends zero to $e_1$ and the $g$-block labels are $0,g^{-1}e_1$. Uniform $g$ makes the latter uniform on $W\setminus\{0\}$, with equal stabilizer multiplicities. Thus rows $24,25$ have exactly the moment in \eqref{eq:power-character}; rows $23,26$ have at most $4^\nu$ times it. Their clock factor is an identical multiplicity, as these words contain only $t,x,z$.

\begin{proof}[Proof of \cref{thm:power-moments}]
Choose
\[
 M=\max\{48,\lceil N^{1/4}\rceil,\lceil(n_0-4)/4\rceil\},\qquad n=4M+4.
\]
The complete space has dimension $D_N=M^2|\operatorname{GL}(n,2)|2^{3N}$, hence
$\log D_N=3N+O_\nu(\sqrt N)$, since $n_0=O_\nu(\log N)$ and $\log|\operatorname{GL}(n,2)|\le n^2$. Every nonzero host residual is at most $3h_M\le24\pi M^{-2}$, whose $2\nu$-moment is $O_\nu(N^{-\nu})$. The character calculation gives the same order for all lamp rows; the marker residual is zero. This proves the theorem for $\nu\ge1$. For $0<\nu<1$, use the $\nu=1$ assignments and scalar concavity:
$\tau((|R|^2)^\nu)\le\tau(|R|^2)^\nu$.
\end{proof}

\begin{corollary}[Explicit rational powers]\label{cor:power-rational}
For integers $1\le a\le b$, one fixed rational BCS game has
\[
 c\eps^{-a/b}2^{-C(\log\log(1/\eps))^2}\le d_G^\to(\eps)
 \le\log D_G(\eps)\le C'\eps^{-a/b}.
\]
\end{corollary}
\begin{proof}
For each normalized source residual replace its zero equation by $R=VS^a$, $S^b=0$, with $V$ unitary and $S$ a positive contraction. At weighted precision $\delta$, the spectral measure $\Tr(\rho1_B(S))$ gives $|S^a|_\rho\le|S^b|_\rho^{a/b}$, without requiring $[S,\rho]=0$. Thus $|R|_\rho\le2\delta^{a/b}$. The compiler's original decoded words have actual mirrors, inverse mirrors and state commutators $O(\sqrt\eps)$, within $e=O(\eps^{a/(2b)})$. The marker gives fixed visibility, so the explicit budget behind \eqref{eq:power-extraction} proves the lower bound on every original mixed state.
For the upper bound take moment order $\nu=b/a$, extend each square-matrix polar partial isometry to a unitary $R=V|R|$, and set $S=|R|^{1/a}$. The first equation is exact and $\|S^b\|_{2,D_N}^2\le C N^{-b/a}$. Flat commutators vanish. Compiler completeness gives loss $O(N^{-b/a})$ with only a constant dimension multiplier and a maximally entangled original state. An integer $N\asymp\eps^{-a/b}$ covers all small losses. Upper losses tend to zero, and the divergent lower bound excludes finite perfection.
\end{proof}

These are finite-moment models, not an operator-norm endpoint theorem: for $u_*=(1,0,\ldots,0)$ all character signs are $-1$, so $X_{u_*}=X_0$, $Z_{u_*}=Z_0$ and the mixed residual has norm one on an exceptional affine block. In contrast its higher finite moments decay. Choosing a uniformly small-bias list gives a separate operator-norm upper estimate $\log D_\eta=O(\eta^{-5/2})$: independent characters obey $\Pr(|c(u)|>\kappa)\le2e^{-N\kappa^2/2}$ by $\cosh t\le e^{t^2/2}$. A union bound gives a list for $N\ge2\kappa^{-2}\ln(4(2^n-1))$, selectable by finite exact search. Take $M\asymp\eta^{-1/2}$, $n=4M+4$, $\kappa=\eta/2$ in the table. This neither proves nor refutes a quadratic operator-norm cost.


The distinction between fixed-word moments and operator-norm control matters in the surrounding approximation literature. Ozawa--R\o rdam--Sato~\cite{ORS14} prove qualitative quasidiagonality for elementary amenable groups, without the dimension rate needed here. Gowers--Hatami~\cite{GH17} study normalized Schatten approximation, and uniform pair-defect stability, while Fournier-Facio--Rangarajan~\cite{FFR23} establish uniform Ulam stability for lamplighters and related groups. Becker--Lubotzky~\cite{BL18} and Arzhantseva--P\u aunescu~\cite{AP14} address stability in their specified Hilbert--Schmidt or Hamming metrics. None of these hypotheses turns a small exceptional fixed-relator block into a small operator norm. No quadratic general filling estimate is assumed from the Houghton literature~\cite{Lee12}; the proof uses the explicit clock regions and their checked affine lift.
